% 第4章 正弦电流电路
\chapter{正弦电流电路}

\section{正弦电流}
随正弦规律变化的电流称为正弦电流或正弦交流

\section{正弦量的相量表示法}
复数

正弦量的相量表示
\[
\dot{A}_m = A_m e^{j\psi} = A_m \angle \psi
\]

正弦量运算与相量运算的对应关系
1.唯一性
2.线性性质
3.微分规则

\section{基尔霍夫的相量形式}

\section{RLC元件上电压与电流的关系}
电阻元件

电感元件
可得电感上电压、电流相量关系
\[
\]

电容元件

\section{阻抗与导纳}
阻抗
即阻抗端口电压、电流有效值（或振幅）之比等于阻抗模；端口电压超前于电流的相位差等于阻抗角$\psi$。

导纳

\section{正弦电流电路的相量分析法}

\section{正弦电流电路的功率}
瞬时功率

平均功率

无功功率
定义无功功率
\[
\]

视在功率

功率因数的提高

复功率

\section{最大功率传输定理}

\section{耦合电感}
互感元件端口特性方程

互感元件的等效电路和去耦合等效

\section{含互感元件的正弦电流电路}
含互感电路的等效化简

\section{理想变压器}